\documentclass[10pt]{report}
\usepackage{anyfontsize}
\AtBeginDocument{\fontsize{8.5}{10.6}\selectfont}

% ===================== ENGINE GUARD =====================
\usepackage{iftex}
\ifXeTeX\else
  \errmessage{This document must be compiled with XeLaTeX.}
\fi

% ===================== PAGE LAYOUT =====================
\usepackage[top=2cm,bottom=2.5cm,left=3cm,right=2cm]{geometry}
\setlength{\headheight}{15pt}

% ===================== CORE PACKAGES =====================
\usepackage{graphicx}
\usepackage{float}
\usepackage{array}
\usepackage{tabularx}
\usepackage{booktabs}
\usepackage{multirow}
\usepackage{threeparttable}
\usepackage{colortbl}
\usepackage{enumitem}
\usepackage{amsmath}
\usepackage{amssymb}
\usepackage{ragged2e}
\usepackage{setspace}

% ===================== VISUAL DESIGN =====================
\usepackage[table]{xcolor}
\usepackage{tcolorbox}
\tcbuselibrary{breakable,skins,theorems}
\usepackage{framed}

\usepackage{fancyhdr}
\usepackage{titlesec}
\usepackage{tocloft}

% ===================== HEADER/FOOTER STYLING =====================
\pagestyle{fancy}
\fancyhf{}
\fancyhead[RE]{\nouppercase{\rightmark}}
\fancyhead[LO]{\nouppercase{\leftmark}}
\fancyhead[LE,RO]{\thepage}
\renewcommand{\headrulewidth}{0.4pt}
\renewcommand{\footrulewidth}{0pt}
\fancypagestyle{plain}{%
  \fancyhf{}
  \fancyfoot[C]{\thepage}
  \renewcommand{\headrulewidth}{0pt}
}

% ===================== PERSIAN (XeLaTeX) =====================
\usepackage{xepersian}

% ============ LATIN FONT SETUP ============
\setlatintextfont[
  Scale=0.95,
  Numbers=Lining,
  Ligatures=TeX
]{TeX Gyre Termes}

% ============ PERSIAN FONT SETUP ============
\settextfont[
  Path=fonts/,
  Extension=.ttf,
  UprightFont=*,
  BoldFont=*Bd,
  ItalicFont=*It,
  BoldItalicFont=*BdIt,
  Script=Arabic,
  Language=Persian
]{XB-Yagut}

\setdigitfont[
  Path=fonts/,
  Extension=.ttf,
  UprightFont=*,
  BoldFont=*Bd,
  ItalicFont=*It,
  BoldItalicFont=*BdIt,
  Numbers=Proportional,
  Script=Arabic,
  Language=Persian
]{XB-Yagut}

% ============ CHAPTER/TITLE FORMATTING ============
\titleformat{\chapter}[display]
  {\normalfont\huge\bfseries\color{PrimaryColor}}
  {\flushright\chaptertitlename\ \thechapter}
  {20pt}
  {\Huge\flushright}
\titlespacing*{\chapter}{0pt}{-20pt}{30pt}

\titleformat{\section}
  {\normalfont\Large\bfseries\color{PrimaryColor}}
  {\thesection}{1em}{}
\titlespacing{\section}{0pt}{12pt}{6pt}

\titleformat{\subsection}
  {\normalfont\large\bfseries\color{SecondaryColor}}
  {\thesubsection}{1em}{}
\titlespacing{\subsection}{0pt}{10pt}{4pt}

% ============ TABLE OF CONTENTS STYLING ============
\renewcommand{\cftchapfont}{\bfseries\color{PrimaryColor}}
\renewcommand{\cftchapleader}{\cftdotfill{\cftdotsep}}
\setlength{\cftbeforechapskip}{6pt}
\renewcommand{\contentsname}{\hfill\Large\bfseries\color{PrimaryColor}فهرست مطالب\hfill}

% ============ ENHANCED MACROS FOR BILINGUAL CONTENT ============
\newcommand{\lat}[1]{\lr{#1}}
\newcommand{\latinterm}[1]{\lr{\textit{#1}}}
\newcommand{\latinbold}[1]{\lr{\textbf{#1}}}
\newcommand{\latinnum}[1]{\lr{#1}}
\newcommand{\technical}[1]{\lr{\textsc{#1}}}

% ============ TABLE COLUMN TYPES FOR RTL ============
\newcolumntype{P}[1]{>{\RaggedRight}p{#1}}
\newcolumntype{H}[1]{>{\bfseries\color{PrimaryColor}\RaggedRight}p{#1}}
\newcolumntype{C}[1]{>{\centering\arraybackslash}p{#1}}

% ============ LINE SPACING ============
\linespread{1.15}

% ============ HYPERREF SETUP ============
\usepackage{hyperref}
\hypersetup{
    colorlinks=true,
    linkcolor=PrimaryColor,
    citecolor=PrimaryColor,
    filecolor=PrimaryColor,
    urlcolor=PrimaryColor,
    pdftitle={آرمان‌های موازی: راهبردهای نیمه‌هادی هند و کره جنوبی},
    pdfauthor={تحلیل تطبیقی سیاست‌گذاری و همکاری‌های بین‌المللی},
    pdfsubject={مطالعه تطبیقی راهبردهای نیمه‌هادی هند و کره جنوبی},
    pdfkeywords={نیمه‌هادی, هند, کره جنوبی, سیاست صنعتی, همکاری بین‌المللی, زنجیره تأمین},
    bookmarksopen=true,
    pdfpagemode=UseOutlines
}

% ===================== PROFESSIONAL COLOR SCHEME =====================
\definecolor{PrimaryColor}{RGB}{44,62,80}
\definecolor{SecondaryColor}{RGB}{90, 70, 70}
\definecolor{IndiaBlue}{RGB}{0, 56, 147}
\definecolor{KoreaRed}{RGB}{205, 45, 45}
\definecolor{LightBackground}{RGB}{248, 248, 252}
\definecolor{IndiaBackground}{RGB}{235, 242, 248}
\definecolor{KoreaBackground}{RGB}{248, 243, 238}
\definecolor{ComparisonBackground}{RGB}{238, 246, 240}
\definecolor{CollaborationBackground}{RGB}{250, 248, 238}
\definecolor{TableHeader}{RGB}{44, 62, 80}
\definecolor{tableRow1}{RGB}{255, 255, 255}
\definecolor{tableRow2}{RGB}{248, 249, 250}

% ===================== ENHANCED TCOLORBOX STYLES WITH CURVATURE =====================
\newtcolorbox{keyconceptbox}[1][]{
  enhanced,
  colback=IndiaBackground,
  colframe=PrimaryColor,
  boxrule=0.75pt,
  arc=4pt,
  fonttitle=\bfseries,
  title={\color{white}مفهوم کلیدی},
  coltitle=white,
  attach boxed title to top right={yshift=-3mm, xshift=-4mm},
  boxed title style={
    colback=PrimaryColor,
    colframe=PrimaryColor,
    coltext=white,
    boxrule=0pt,
    arc=3pt
  },
  breakable,
  left=8pt,
  right=8pt,
  top=10pt,
  bottom=8pt,
  #1
}

\newtcolorbox{researchbox}[1][]{
  enhanced,
  colback=LightBackground,
  colframe=PrimaryColor,
  boxrule=0.75pt,
  arc=4pt,
  fonttitle=\bfseries,
  title={\color{PrimaryColor}یادداشت پژوهشی},
  breakable,
  left=8pt,
  right=8pt,
  top=10pt,
  bottom=8pt,
  #1
}

\newtcolorbox{indiabox}[2][]{
  enhanced,
  colback=IndiaBackground,
  colframe=IndiaBlue,
  boxrule=0.75pt,
  arc=4pt,
  fonttitle=\bfseries,
  title=#2,
  attach boxed title to top left={yshift=-3mm, xshift=4mm},
  boxed title style={
    colback=IndiaBlue,
    colframe=IndiaBlue,
    coltext=white,
    boxrule=0pt,
    arc=3pt
  },
  breakable,
  left=8pt,
  right=8pt,
  top=10pt,
  bottom=8pt,
  #1
}

\newtcolorbox{koreabox}[2][]{
  enhanced,
  colback=KoreaBackground,
  colframe=KoreaRed,
  boxrule=0.75pt,
  arc=4pt,
  fonttitle=\bfseries,
  title=#2,
  attach boxed title to top left={yshift=-3mm, xshift=4mm},
  boxed title style={
    colback=KoreaRed,
    colframe=KoreaRed,
    coltext=white,
    boxrule=0pt,
    arc=3pt
  },
  breakable,
  left=8pt,
  right=8pt,
  top=10pt,
  bottom=8pt,
  #1
}

\newtcolorbox{comparisonbox}[2][]{
  enhanced,
  colback=ComparisonBackground,
  colframe=PrimaryColor,
  boxrule=0.75pt,
  arc=4pt,
  fonttitle=\bfseries,
  title=#2,
  attach boxed title to top left={yshift=-3mm, xshift=4mm},
  boxed title style={
    colback=PrimaryColor,
    colframe=PrimaryColor,
    coltext=white,
    boxrule=0pt,
    arc=3pt
  },
  breakable,
  left=8pt,
  right=8pt,
  top=10pt,
  bottom=8pt,
  #1
}

\newtcolorbox{collaborationbox}[2][]{
  enhanced,
  colback=CollaborationBackground,
  colframe=SecondaryColor,
  boxrule=0.75pt,
  arc=4pt,
  fonttitle=\bfseries,
  title=#2,
  attach boxed title to top left={yshift=-3mm, xshift=4mm},
  boxed title style={
    colback=SecondaryColor,
    colframe=SecondaryColor,
    coltext=white,
    boxrule=0pt,
    arc=3pt
  },
  breakable,
  left=8pt,
  right=8pt,
  top=10pt,
  bottom=8pt,
  #1
}

\newtcolorbox{summarybox}[2][]{
  enhanced,
  colback=LightBackground,
  colframe=SecondaryColor,
  boxrule=0.75pt,
  arc=4pt,
  fonttitle=\bfseries,
  title=#2,
  attach boxed title to top left={yshift=-3mm, xshift=4mm},
  boxed title style={
    colback=SecondaryColor,
    colframe=SecondaryColor,
    coltext=white,
    boxrule=0pt,
    arc=3pt
  },
  breakable,
  left=8pt,
  right=8pt,
  top=10pt,
  bottom=8pt,
  #1
}

% ===================== CUSTOM COMMANDS =====================
\newcommand{\indiatext}[1]{\textcolor{IndiaBlue}{\textbf{#1}}}
\newcommand{\koreatext}[1]{\textcolor{KoreaRed}{\textbf{#1}}}

% ===================== DOCUMENT STARTS =====================
\begin{document}

% ===================== صفحه عنوان =====================
\thispagestyle{empty}
\begin{center}
\vspace*{2cm}

{\Huge\bfseries\color{PrimaryColor} آرمان‌های موازی: راهبردهای نیمه‌هادی هند و کره جنوبی}

\vspace{0.6cm}
{\color{gray}\rule{0.7\textwidth}{0.75pt}}
\vspace{0.6cm}

{\Large\bfseries تحلیلی تطبیقی از سیاست‌گذاری، توسعه زیست‌بوم، و شبکه‌های همکاری جهانی}

\vspace{1.5cm}

\begin{researchbox}
\centering
\textbf{این گزارش تحلیلی موازی از راهبردهای نیمه‌هادی دنبال‌شده توسط هند و کره جنوبی ارائه می‌دهد. این گزارش به بررسی سیاست‌های بنیادین، توسعه زیست‌بوم داخلی، و شبکه‌های متمایز مشارکت‌های بین‌المللی می‌پردازد که هر کشور برای تضمین موقعیت خود در زنجیره ارزش جهانی در حال پرورش آنهاست. این تحلیل بر ساخت‌وساز ماموریت‌محور هند از یک پایه طراحی قوی و تلاش کره جنوبی برای حفظ رهبری از طریق سرمایه‌گذاری و تحقیق و توسعه کلان تأکید دارد.}
\end{researchbox}

\vfill
{\small\color{gray} نسخه سند: ۱.۰ \\ تاریخ: \today}
\end{center}

% ===================== فهرست مطالب =====================
\clearpage
\tableofcontents

% ===================== فصل 1: مقدمه =====================
\clearpage
\chapter{مقدمه: نیمه‌هادی‌ها به مثابه یک ضرورت ژئواستراتژیک}

\begin{keyconceptbox}
صنعت جهانی نیمه‌هادی‌ها که پایه پیشرفت‌ها از هوش مصنوعی تا سیستم‌های دفاعی است، در حال بازترازیابی قابل توجهی است. با محرک‌هایی مانند آسیب‌پذیری‌های زنجیره تأمین که در طول همه‌گیری کووید-۱۹ آشکار شد و رقابت فناورانه فزاینده آمریکا و چین، کشورها به طور تهاجمی در پی راهبردهایی برای انعطاف‌پذیری و خودکفایی بیشتر هستند.
\end{keyconceptbox}

این گزارش تحلیلی تطبیقی از دو بازیگر کلیدی در این بازترازیابی ارائه می‌دهد: \indiatext{هند}، یک قطب بلندپرواز در حال ظهور که در حال ساخت یک زیست‌بوم جامع از پایه است، و \koreatext{کره جنوبی}، یک قدرت تأسیس‌شده که برای حفظ برتری خود در حافظه و ساخت پیشرفته می‌کوشد.

این تحلیل به گونه‌ای ساختار یافته است که هر دو کشور را با عمق یکسان مورد بررسی قرار دهد. نخست، راهبردهای ملی محوری و چارچوب‌های سیاستی آنها را ترسیم خواهد کرد. دوم، ترکیب و رشد زیست‌بوم‌های داخلی نیمه‌هادی آنها را ترسیم می‌کند. بخش نهایی و گسترده‌ترین بخش، \textbf{شبکه‌های همکاری بین‌المللی} مربوطه آنها را تشریح می‌کند و به جزئیات شرکای کلیدی، ماهیت همکاری و اهداف استراتژیک زیربنایی می‌پردازد.

% ===================== فصل 2: راهبردهای ملی =====================
\chapter{راهبردها و چارچوب‌های سیاستی ملی}

\section{ماموریت نیمه‌هادی هند: از بنیان تا زیست‌بوم}

\begin{indiabox}{ISM 1.0: گذاشتن بنیان}
ماموریت نیمه‌هادی هند (ISM 1.0) که در دسامبر ۲۰۲۱ راه‌اندازی شد، یک چارچوب مشوق ۱۰ میلیارد دلاری بود که برای جذب سرمایه‌گذاری‌های کلان در حوزه‌های ساخت (فاب)، مونتاژ، آزمایش، نشانه‌گذاری و بسته‌بندی (ATMP/OSAT) و طراحی تراشه طراحی شده بود. موفقیت اولیه آن در بذرافشانی زیست‌بوم بوده است، به طوری که تا دسامبر ۲۰۲۵، ۱۰ پروژه اصلی در شش ایالت تأیید شده است که نشان‌دهنده سرمایه‌گذاری کلی حدود ۱.۶۰ لاخ کرور روپیه است.
\end{indiabox}

\subsection{ISM 2.0: چرخش استراتژیک به سمت طراحی و عمق}
ISM 2.0 که در بودجه اتحادیه ۲۰۲۷-۲۰۲۶ اعلام شد، نشان‌دهنده تکاملی استراتژیک با پیش‌بینی اولیه ۱۰۰۰ کرور روپیه است. تمرکز آن از جذب سرمایه به سمت تعمیق قابلیت فناورانه تغییر می‌کند:

\begin{comparisonbox}{تمرکز‌های اصلی ISM 2.0}
\begin{itemize}
    \item \textbf{رشد طراحی‌محور:} اولویت دادن به توسعه مالکیت فکری (IP) «تمام‌پشته» هندی و حمایت از استارت‌آپ‌های فاب‌لس. یک پروژه توسعه مشترک برای یک تراشه ۲ نانومتری با کوالکام نمونه‌ای از این بلندپروازی است.
    \item \textbf{استحکام‌بخشیدن به زیست‌بوم:} هدف گذاری برای ایجاد تولید داخلی تجهیزات، مواد و گازهای ویژه نیمه‌هادی به منظور کاهش وابستگی به واردات.
    \item \textbf{خط لوله استعدادها:} گسترش ابتکارهایی مانند برنامه «از تراشه تا استارت‌آپ» که دسترسی به ابزار را در ۳۹۷ مؤسسه دانشگاهی فراهم می‌کند، برای ساخت نیروی کار متخصص.
\end{itemize}
\end{comparisonbox}

نقشه راه بلندمدت صراحت دارد: دستیابی به قابلیت‌های ساخت ۳ نانومتر و ۲ نانومتر و قرار دادن هند در میان کشورهای برتر جهانی نیمه‌هادی تا سال ۲۰۳۵.

\section{راهبرد کره‌ای: دفاع از رهبری از طریق سرمایه‌گذاری}

\begin{koreabox}{ابتکار کمربند نیمه‌هادی کره}
در پاسخ به فشارهای زنجیره تأمین جهانی، کره جنوبی در سال ۲۰۲۱ راهبرد ۴۵۰ میلیارد دلاری «کمربند نیمه‌هادی کره» را راه‌اندازی کرد. این طرح ملی با هدف تحکیم یک خوشه ساخت منطقه‌ای و تضمین رهبری در فناوری‌های نسل بعدی مانند حافظه پهنای باند بالا (HBM) برای هوش مصنوعی است.
\end{koreabox}

\subsubsection{تقویت سیاست: قانون تراشه‌های کره}
قانون تراشه‌های کره که در فوریه ۲۰۲۵ به تصویب رسید، مشوق‌های مالیاتی قابل توجهی برای تحقیق و توسعه داخلی و سرمایه‌گذاری‌های تأسیساتی فراهم می‌کند و محیط مالی مساعدی برای غول‌هایی مانند سامسونگ و اس‌کهاینیکس ایجاد می‌کند تا ظرفیت ساخت پیشرفته خود را در داخل کشور گسترش دهند.

\subsubsection{وابستگی اقتصادی و چشم‌انداز ۲۰۲۶}
\begin{researchbox}
نیمه‌هادی‌ها سنگ بنای اقتصاد کره جنوبی هستند و صادرات آن را در سال ۲۰۲۵ از مرز ۷۰۰ میلیارد دلار گذراندند. برای سال ۲۰۲۶، دولت رشد اقتصادی ۲ درصدی را پیش‌بینی می‌کند که شدیداً متکی به ادامۀ «رونق تراشه» است و پیش‌بینی می‌شود هزینه‌های تجهیزات داخلی ۲۷.۲ درصد افزایش یابد و به نزدیک ۳۰ میلیارد دلار برسد، زیرا شرکت‌ها تولید خود را برای حافظه هوش مصنوعی افزایش می‌دهند.
\end{researchbox}

% ===================== فصل 3: زیست‌بوم‌های داخلی =====================
\chapter{ترکیب زیست‌بوم داخلی}

\section{زیست‌بوم نوپا اما به سرعت در حال گسترش هند}
زیست‌بوم هند با نقاط ورود استراتژیک و بنیادی قوی در طراحی مشخص می‌شود.

\begin{table}[H]
    \centering
    \caption{ستون‌های کلیدی زیست‌بوم نیمه‌هادی هند}
    \label{tab:india_ecosystem}
    \begin{tabular}{H{0.25\linewidth} P{0.7\linewidth}}
        \toprule
        \rowcolor{TableHeader}
        \textcolor{white}{\textbf{ستون}} & \textcolor{white}{\textbf{شرح و وضعیت}} \\
        \midrule
        \rowcolor{tableRow1}
        \textbf{ساخت} & اولین کارخانه تولید تجاری از طریق سرمایه‌گذاری مشترک تاتا-PSMC در گجرات در حال انجام است. کارخانه‌های نیمه‌هادی ترکیبی (مانند SiC) نیز تأیید شده‌اند. تمرکز در ابتدا بر روی گره‌های بالغ (۶۵-۲۸ نانومتر) و میانی باقی می‌ماند. \\
        \rowcolor{tableRow2}
        \textbf{ATMP/OSAT} & یک حوزه موفقیت کلیدی. تأسیسات میکرون در گجرات عملیاتی است. چندین شرکت هندی (تاتا، کینز) و سرمایه‌گذاری‌های مشترک (CG Power با رنسان) در حال ایجاد ظرفیت بسته‌بندی هستند. \\
        \rowcolor{tableRow1}
        \textbf{طراحی تراشه} & نقطه قوت اصلی، با اشتغال‌زایی تخمینی ۲۰ درصد از استعدادهای طراحی جهانی. طرح DLI از ۲۴ استارت‌آپ حمایت می‌کند که ۱۶ نمونه اولیه (tape-out) از جمله در گره ۱۲ نانومتر را تکمیل کرده‌اند. \\
        \rowcolor{tableRow2}
        \textbf{مالکیت فکری بومی} & توسعه معماری‌های پردازنده بومی (SHAKTI, DHRUV64) بر پایه RISC-V متن‌باز، با هدف حاکمیت فناورانه. \\
        \bottomrule
    \end{tabular}
\end{table}

\section{زیست‌بوم بالغ و متمرکز کره جنوبی}
زیست‌بوم کره جنوبی توسط قهرمانان یکپارچه عمودی تسلط دارد و در حافظه در سطح جهانی پیشرو است.

\begin{table}[H]
    \centering
    \caption{ستون‌های کلیدی زیست‌بوم نیمه‌هادی کره جنوبی}
    \label{tab:korea_ecosystem}
    \begin{tabular}{H{0.25\linewidth} P{0.7\linewidth}}
        \toprule
        \rowcolor{TableHeader}
        \textcolor{white}{\textbf{ستون}} & \textcolor{white}{\textbf{شرح و وضعیت}} \\
        \midrule
        \rowcolor{tableRow1}
        \textbf{ساخت حافظه} & سلطۀ پیشرو در جهان. بیش از ۷۰ درصد از DRAM جهانی و بیش از نیمی از NAND فلش را تولید می‌کند، عمدتاً از طریق سامسونگ و اس‌کهاینیکس. تأمین‌کننده حیاتی HBM برای هوش مصنوعی. \\
        \rowcolor{tableRow2}
        \textbf{کارخانه و منطق} & سامسونگ یک بازیکن برتر در حوزه کارخانه (foundry) است که در مرز گره‌های پیشرفته (مانند ۳ نانومتر) رقابت می‌کند. یک سازنده قراردادی کلیدی برای تراشه‌های منطقی جهانی. \\
        \rowcolor{tableRow1}
        \textbf{تجهیزات و مواد} & زنجیره تأمین داخلی قوی، اگرچه متکی به واردات کلیدی (مانند لیتوگرافی EUV). هزینه تجهیزات یک شاخص رشد عمده است. \\
        \rowcolor{tableRow2}
        \textbf{تحقیق و توسعه و نوآوری} & سرمایه‌گذاری سنگین در فناوری نسل بعدی (مانند تراشه‌های هوش مصنوعی، بسته‌بندی سه‌بعدی، طراحی چیپلت). انجمن‌های صنعتی مانند SEMICON Korea 2026 دستورکارهای جهانی را تعیین می‌کنند. \\
        \bottomrule
    \end{tabular}
\end{table}

% ===================== فصل 4: شبکه‌های همکاری بین‌المللی =====================
\chapter{شبکه‌های همکاری بین‌المللی}
این بخش مشارکت‌های دوجانبه و چندجانبه اولیه دنبال‌شده توسط هند و کره جنوبی را ترسیم می‌کند و ماهیت مکمل و استراتژیک آنها را برجسته می‌سازد.

\section{مشارکت استراتژیک هند و کره جنوبی}
\begin{collaborationbox}{مکمل‌های استراتژیک}
این رابطه دوجانبه سنگ بنای راهبردهای هر دو ملت است که بر پایه مکمل‌های واضحی بنا شده است: سرمایه، فناوری و قدرت ساخت کره با پتانسیل بازار، استعداد طراحی و محیط پایدار سرمایه‌گذاری هند دیدار می‌کند.
\end{collaborationbox}

\begin{table}[H]
    \centering
    \caption{ماتریس مشارکت نیمه‌هادی هند و کره جنوبی}
    \label{tab:india_korea_matrix}
    \begin{tabular}{P{0.3\linewidth} P{0.32\linewidth} P{0.3\linewidth}}
        \toprule
        \rowcolor{TableHeader}
        \textcolor{white}{\textbf{شکل همکاری}} & \textcolor{white}{\textbf{منافع و دستاوردهای هند}} & \textcolor{white}{\textbf{منافع و دستاوردهای کره جنوبی}} \\
        \midrule
        \rowcolor{tableRow1}
        \textbf{سرمایه‌گذاری و ساخت} & جذب سرمایه و تأسیسات پیشرفته (مانند ارزیابی اس‌کهاینیکس برای یک کارخانه بسته‌بندی) برای ساخت ظرفیت داخلی. & متنوع‌سازی پایگاه‌های تولید دور از مخاطرات ژئوپلیتیکی، دسترسی به مشوق‌ها و تضمین یک بازار رشد بلندمدت. \\
        \rowcolor{tableRow2}
        \textbf{فناوری و تحقیق و توسعه} & کسب تخصص در حوزه‌های پیشرفته مانند تراشه‌های هوش مصنوعی و بسته‌بندی سه‌بعدی از طریق مراکز تحقیقاتی مشترک. & بهره‌گیری از استعداد طراحی هندی، همکاری بر روی راه‌حل‌های نوآورانه و شکل‌دهی به استانداردهای آینده. \\
        \rowcolor{tableRow1}
        \textbf{یکپارچگی زنجیره تأمین} & ساخت شبکه‌های تأمین‌کننده (مانند تفاهم‌نامه ودانتا با ۲۰ شرکت شیشه نمایشگر کره‌ای) برای عمق بخشیدن به زیست‌بوم. & ایجاد یک زنجیره تأمین انعطاف‌پذیر به عنوان جایگزین چین برای مواد و قطعات. \\
        \rowcolor{tableRow2}
        \textbf{گفت‌وگوی نهادی} & انجمن‌هایی مانند مجمع اقتصاد آینده کره-هند همسویی را تقویت کرده و فرصت‌های جدید را شناسایی می‌کنند. & تضمین پیش‌بینی‌پذیری سیاست و تقویت روابط استراتژیک در هندو-آرام. \\
        \bottomrule
    \end{tabular}
\end{table}

\section{شبکه همکاری هند: ساخت یک زیست‌بوم جهانی}
راهبرد مشارکت هند چندجهته است و هدف آن کسب فناوری، جذب سرمایه‌گذاری و ادغام در زنجیره‌های تأمین قابل اعتماد است.

\begin{table}[H]
    \centering
    \caption{مشارکت‌های کلیدی نیمه‌هادی هند}
    \label{tab:india_partners}
    \begin{tabular}{P{0.22\linewidth} P{0.35\linewidth} P{0.35\linewidth}}
        \toprule
        \rowcolor{TableHeader}
        \textcolor{white}{\textbf{کشور/منطقه شریک}} & \textcolor{white}{\textbf{ماهیت همکاری}} & \textcolor{white}{\textbf{هدف استراتژیک برای هند}} \\
        \midrule
        \rowcolor{tableRow1}
        \textbf{ایالات متحده} & همسویی استراتژیک، سرمایه‌گذاری از شرکت‌هایی مانند میکرون (ATMP) و همکاری طراحی (مانند با کوالکام). & ادغام در یک اتحاد فناوری «مورد اعتماد» غربی، جذب فناوری و سرمایه پیشرو. \\
        \rowcolor{tableRow2}
        \textbf{تایوان} & انتقال فناوری حیاتی و تخصص عملیاتی از طریق سرمایه‌گذاری‌های مشترک (مانند تاتا با PSMC برای کارخانه). & راه‌اندازی قابلیت‌های ساخت پیشرفته با یک شریک اثبات‌شده. \\
        \rowcolor{tableRow1}
        \textbf{ژاپن} & مشارکت در بسته‌بندی (مانند CG Power با رنسانس) و تأمین مواد/تجهیزات. & دسترسی به تخصص خاص، متنوع‌سازی مشارکت‌ها فراتر از تایوان/کره و ساخت یک زنجیره تأمین انعطاف‌پذیر. \\
        \rowcolor{tableRow2}
        \textbf{اتحادیه اروپا} & همکاری‌های در حال ظهور در تجهیزات و تراشه‌های خودرو. & دسترسی به فناوری پیشرفته و پاسخگویی به بخش ساخت خودرو/خودروی الکتریکی در حال رشد در هند. \\
        \bottomrule
    \end{tabular}
\end{table}

\section{شبکه همکاری کره جنوبی: تضمین بازارها و نوآوری}
راهبرد کره جنوبی بر تضمین بازار صادراتی خود، دسترسی به نوآوری پیشرو و کاهش ریسک ژئوپلیتیکی متمرکز است.

\begin{table}[H]
    \centering
    \caption{مشارکت‌های کلیدی نیمه‌هادی کره جنوبی}
    \label{tab:korea_partners}
    \begin{tabular}{P{0.22\linewidth} P{0.35\linewidth} P{0.35\linewidth}}
        \toprule
        \rowcolor{TableHeader}
        \textcolor{white}{\textbf{کشور/منطقه شریک}} & \textcolor{white}{\textbf{ماهیت همکاری}} & \textcolor{white}{\textbf{هدف استراتژیک برای کره جنوبی}} \\
        \midrule
        \rowcolor{tableRow1}
        \textbf{ایالات متحده} & روابط عمیق با مشتریان (مانند تأمین HBM به انویدیا)، و حرکت در میان محدودیت‌های فناورانه تحت رهبری آمریکا. & حفظ دسترسی به بازار اصلی تراشه‌های هوش مصنوعی و تجهیزات/نرم‌افزارهای حیاتی، مدیریت انطباق ژئوپلیتیکی. \\
        \rowcolor{tableRow2}
        \textbf{هند} & همان‌طور که در جدول \ref{tab:india_korea_matrix} جزئیات داده شد: سرمایه‌گذاری، دسترسی به استعدادها، متنوع‌سازی زنجیره تأمین. & متنوع‌سازی تولید، تضمین بازار رشد آینده و ایجاد عمق استراتژیک در آسیا. \\
        \rowcolor{tableRow1}
        \textbf{ژاپن} & رابطه پیچیده شامل رقابت و همکاری در مواد/تجهیزات. شریک بالقوه در یک چارچوب سه‌جانبه گسترده‌تر با هند. & تضمین تأمین مواد حیاتی و کاوش چارچوب‌های همکاری برای کاهش ریسک‌های سیستمی. \\
        \rowcolor{tableRow2}
        \textbf{اتحادیه اروپا / هلند} & رابطه تأمین‌کننده کلیدی برای تجهیزات ساخت حیاتی (مانند ASML). & تضمین دسترسی بی‌وقفه به پیشرفته‌ترین ابزارهای ساخت. \\
        \bottomrule
    \end{tabular}
\end{table}

\section{پتانسیل یک سه‌جانبه هند-ژاپن-کره جنوبی}
\begin{keyconceptbox}
تحلیلگران پیشنهاد می‌کنند که یک چارچوب سه‌جانبه رسمی می‌تواند به طور قدرتمندی نقاط قوت هر سه را هم‌افزایی کند: \textbf{استعداد طراحی و بازار هند}، \textbf{تسلط ژاپن بر تجهیزات و مواد}، و \textbf{ساخت در کلاس جهانی کره جنوبی}. چنین اتحادی با هدف ایجاد یک زیست‌بوم نیمه‌هادی «مقاوم در برابر چین» و خودکفاتر در داخل منطقه هندو-آرام خواهد بود، و وابستگی بیش از حد به هر جغرافیای واحد را کاهش داده و خودمختاری استراتژیک جمعی را افزایش خواهد داد.
\end{keyconceptbox}

% ===================== فصل 5: نتیجه‌گیری =====================
\chapter{نتیجه‌گیری: مسیرهای همگرا در جهانی در حال تکه‌تکه شدن}

\begin{summarybox}{نتیجه‌گیری کلیدی}
هند و کره جنوبی، اگرچه در مراحل مختلف توسعه نیمه‌هادی قرار دارند، در مسیرهای استراتژیک همگرایی قرار گرفته‌اند که توسط ضرورت‌های یکسانی هدایت می‌شوند: حاکمیت فناورانه، انعطاف‌پذیری زنجیره تأمین و امنیت اقتصادی. مشارکت دوجانبه آنها عنصری عمل‌گرایانه و به طور فزاینده‌ای مرکزی از هر دو راهبرد است که مزایای متقابل واضحی ارائه می‌دهد.
\end{summarybox}

در نهایت، موفقیت آنها نه تنها به سیاست و سرمایه‌گذاری داخلی، بلکه به توانایی آنها در حرکت و شکل‌دهی به شبکه پیچیده همکاری‌های بین‌المللی در یک صنعت با بار ژئوپلیتیکی بالا بستگی خواهد داشت. ظهور چارچوب‌های همکاری جدید، که به طور بالقوه ژاپن را نیز شامل می‌شود، می‌تواند به طور قابل توجهی چشمانداز نیمه‌هادی منطقه‌ای و جهانی را در دهه آینده تغییر دهد.

% ===================== منابع =====================
\section*{منابع}
\begin{enumerate}[rightmargin=1.5cm]
    \item دفتر اطلاعات رسانی همگانی. (۷ فوریه ۲۰۲۶). \textit{ماموریت نیمه‌هادی هند ۲.۰}. دولت هند.
    \item KED Global. (۹ ژانویه ۲۰۲۶). \textit{کره جنوبی شرط خود را بر رونق صادرات تراشه برای افزایش رشد اقتصادی ۲۰۲۶ گذاشته است}.
    \item یاداو، ک. ک. (۱۱ ژوئیه ۲۰۲۵). \textit{مشارکت نیمه‌هادی هند و کره جنوبی: هم‌افزایی استراتژیک در میان رقابت آمریکا و چین}. شورای هند امور جهانی.
    \item تریپاتی، پ. (۱۶ اکتبر ۲۰۲۵). \textit{نیمه‌هادی‌ها به عنوان جرقه‌ای برای یک سه‌جانبه هند-ژاپن-کره جنوبی}. بنیاد پژوهشی آبزرور.
    \item SEMI. (۱۵ ژانویه ۲۰۲۶). \textit{SEMICON Korea 2026 سخنرانی‌های کلیدی و کنفرانس‌ها را تقویت می‌کند}.
    \item The Economic Times. (۲۰۲۶). \textit{بودجه اتحادیه ۲۰۲۶: وزیر دارایی از ماموریت نیمه‌هادی هند ۲.۰ رونمایی کرد}.
    \item Whalesbook. (۲۰۲۶). \textit{بلندپروازی تراشه هند: ISM 2.0 طراحی‌محور قطب جهانی را نشانه رفته است}.
    \item NGS India. \textit{روابط هند و کره جنوبی با مشارکت فناوری‌محور قوی‌تر می‌شود}.
\end{enumerate}

\end{document}